\section{Результаты}

В результате получено интерактивное приложение, отображающее усечённый правильный тетраэдр. В обычном режиме различимы треугольные и шестиугольные грани, их контуры и изменение яркости при повороте. Каркасный режим показывает полный набор рёбер и позволяет проследить связи между вершинами.

Проверка геометрии выполнена отдельной программой \path{tetrahedron_test.cpp}. Она проверяет количество вершин, рёбер и граней, наличие четырёх граней каждого типа, одинаковую длину рёбер, плоскостность и правильность граней, выпуклость их контуров и направление нормалей. Дополнительно проверены принадлежность каждого ребра двум граням, наличие трёх рёбер в каждой вершине и сохранение расстояния до центра после поворота. Все проверки завершились успешно.

Проект успешно собран в конфигурации Debug средствами CMake и Ninja. При запуске после исправления записи основного буфера команд сообщения об ошибках Vulkan не появились в наблюдавшемся сеансе. Предупреждения компилятора в сторонней библиотеке VulkanMemoryAllocator не препятствовали сборке.

Работа обеих проекций и каркасного режима проверена на запущенном приложении. При смене перспективной проекции на ортографическую меняется вид фигуры, но сохраняются её геометрия и текущая ориентация. Ползунки и управление мышью позволяют изменять вид без перезапуска. Отображение пересчитывается в основном цикле программы; количественное измерение частоты кадров не проводилось.

Реализация использует все три обязательные библиотеки. Требование динамической смены проекции выполнено переключателем в интерфейсе. Взаимодействие с объектом обеспечено поворотом, масштабом отображения и режимами визуализации. Изменение масштаба отображения не является независимым растяжением самой модели по координатным осям.

Обнаруженное при первоначальном запуске падение было связано с отправкой незаписанного буфера команд и отсутствием подготовки изображения к следующему проходу рисования. В функцию render добавлены запись команд и основной проход очистки. Кроме того, исправлено условие, которое ошибочно выводило сообщение о неудачном показе кадра при успешном результате.

\input{\ReportRoot/sections/results/01-screenshots}
